\scanpage{130}
% archive foliation: 129
écrire $V=\mathbb V(\omega_\sigma)$, $\omega_\sigma$ étant le fibré conormal le long de la section $\sigma$, qui est aussi isom. au fibré conormal de $X_X$ le long de la section diagonale, donc
\[
V=\mathbb V(\underbrace{\Omega^1_{X/\mathbb C}}_{\omega_{X/\mathbb C}})=\mathbb W(\underbrace{t_{X/\mathbb C}}_{\text{fibré tangent}}).
\]
Donc $\widetilde X$ s'identifie au fibré tangent à $X$, $\sigma$ à la section nulle, $P_X$ à la clôture projective de $\widetilde X=\mathbb W(t_{X/\mathbb C})$. La connexion \deleted{intégrable}\inserted{plate} canonique de $P_X$ est associée à une connexion de $\widetilde X$ \uline{considéré comme fibré en droites affines}, i.e. une connex. sur le fibré principal associé de groupe $\operatorname{Aff}(1,\mathbb C)$ (\unsure{le sous-}$\mathbb G_m(\mathbb C)$ de $\widetilde X=\mathbb W(t_{X/\mathbb C})$ privilégié, est contenu dans $\mathrm{GL}(1,\mathbb C)$).

Il y \deleted{aurait lieu}\inserted{aurait à} construire cette connexion, et prouver (ou l'infirmer !) qu'elle est en fait \uline{algébrique}. (Ce qui ne va pas de soi si $X$ non complète, i.e. $X=\mathbb E^1_{\mathbb C}$ ou $X=\mathbb G_{m\mathbb C}$ (cas compact réglé par \unsure{Gaga!|GAGA}). Or dans le cas $X\simeq \mathbb E^1_{\mathbb C}$, on a encore $\widetilde X=X\times X$, $P_X=P(X)\times X$, grâce à $\widetilde X\subset X$, $p:\widetilde X\longrightarrow X_X$ est l'identité (donc $\sigma=\Delta_X$), tout est clair (et $V$ est algébrique --- la connexion dans le cadre algébrique : \unsure{Transport !})
\margin{oui elle est algébrique cf plus bas}

Dans le cas de $X\deleted{\unsure{\simeq}}\simeq\mathbb C^*\simeq\mathbb G_{m\mathbb C}(\mathbb C)$
